\section{Evaluation direction: sequential attack introduction}
\textbf{Design status: not yet an executed or fixed protocol.} The user clarified that reproducing raw collection timestamps is not mandatory, and that comparing performance immediately before and after acquiring a few labels is not a required objective. Sequential introduction is an evaluation axis for scarce-label attack detection; online context updating is a separate design choice.

A conceptual evaluation starts with benign traffic and some attacks, introduces additional attacks in stages, and observes class-specific performance as the mixture changes. An attack first appearing in the evaluation stream may already have a few training examples. It must not be called unseen or out-of-distribution solely because it appears later in the test sequence. Attack order, prevalence, available training labels and context update policy must be explicitly documented and controlled. With a fixed context and the same samples, reordering evaluation alone leaves individual predictions unchanged. Such a scenario measures performance under stage-specific mixtures; it is not evidence of adaptation without an update mechanism. No fixed number of stages, label quota or update interval is adopted in this revision.

\paragraph{Comparisons.} Use the same labelled information and evaluation rows for XGBoost, Global TabPFN, and the proposed context composition. Match total stored context, including repeated anchors, when isolating composition or expert-count effects. Report additional data used to train embeddings, optimise contexts or fit scorer/verifier policies. Introduce component ablations after a useful configuration is identified on validation data.

\paragraph{Measurements.} Report per-attack precision, recall and F1, benign false alarms, retained performance on existing attacks, and context construction/model fitting/inference time and memory. Benign false-alarm rate is the fraction of true benign examples predicted as attacks; it is not the fraction of attacks predicted benign. Report support and prevalence in each stage. Stage macro-F1 is secondary and cannot be compared without accounting for differing class sets. Avoid selecting a preferred attack order or model setting on the final test outcomes.

\paragraph{Optional updates and OOD extension.} If online updates are later introduced, prediction must precede label availability, and only labels available at that moment may affect future predictions. A globally chronological split is needed when claiming natural-time deployment performance. Later OOD evaluation should separately address distribution shifts or attacks absent from training, distinguishing recognition before labels from adaptation after labels. Few-label classification, sequential arrival and zero-label novelty detection are distinct problems.
